\section*{Bab \ref{ch:b2:diffcalc} --- Kalkulus Diferensial}

\begin{solution}{exo:b2:diffcalc:1}
$J_f = \begin{pmatrix} 2x & -2y\\ 2y & 2x\end{pmatrix}$ dengan $\det J_f
= 4(x^2 + y^2)$: jadi terbalikkan jauh dari titik asalnya. (Sebab $f$ ini adalah $z
\mapsto z^2$ dalam samaran kompleks.)

$J_\Phi = \begin{pmatrix} \cos\theta & -r\sin\theta & 0\\
\sin\theta & r\cos\theta & 0\\ 0 & 0 & 1\end{pmatrix}$ dengan $\det = r$:
jadi terbalikkan untuk $r \neq 0$ (yakni koordinat silinder).
\end{solution}

\begin{solution}{exo:b2:diffcalc:2}
Untuk $v = (a, b) \neq 0$: $\frac{f(tv) - 0}{t} = \frac{t^3a^3}{t\cdot
t^2(a^2+b^2)} = \frac{a^3}{a^2 + b^2}$: jadi setiap turunan berarahnya
ada, dengan nilai $D_v = \frac{a^3}{a^2+b^2}$. Tetapi $v
\mapsto D_v$ tidak linear (sebab $D_{(1,0)} = 1$, $D_{(0,1)} = 0$,
$D_{(1,1)} = \frac12 \neq 1$): jadi tak ada pemetaan linear yang dapat menghasilkan nilai
itu, sehingga $f$ tidak \omterm{def:b2:diffcalc:differential}{terdiferensialkan} di $0$ (sebab \omterm{def:b2:diffcalc:differential}{diferensialnya}
mestilah $v \mapsto D_v$).
\end{solution}

\begin{solution}{exo:b2:diffcalc:3}
$f = x^3 + y^3 - 3xy$: titik kritisnya $(0,0)$ dan $(1,1)$ (lewat
perhitungan Tahun ke-1). Hessiannya: $H = \begin{pmatrix} 6x & -3\\ -3 &
6y\end{pmatrix}$. Di $(0,0)$: tanda eigennya bercampur (sebab $\det = -9 < 0$),
jadi pelana. Di $(1,1)$: $\det = 27 > 0$ dan tracenya $> 0$, jadi definit
positif, sehingga minimum lokal yang tegas --- yang kini dibenarkan
\cref{thm:b2:diffcalc:extrema}, bukan sekadar dititahkan.

$g = x^4 + y^4 - 2(x-y)^2$: $\nabla g = (4x^3 - 4(x - y),\; 4y^3 +
4(x-y))$; jadi titik kritisnya $(0,0)$, $(\sqrt2, -\sqrt2)$,
$(-\sqrt2, \sqrt2)$ (lewat Tahun ke-1). Di $(\pm\sqrt2, \mp\sqrt2)$: $H =
\begin{pmatrix} 12\cdot2 - 4 & 4\\ 4 & 20\end{pmatrix} =
\begin{pmatrix} 20 & 4\\ 4 & 20\end{pmatrix}$: jadi definit positif
(sebab diagonalnya dominan; dan \omterm{def:b2:reduction:eigen}{nilai eigennya} $24, 16$): sehingga minimum lokal yang tegas.
Di $(0,0)$: $H = \begin{pmatrix} -4 & 4\\ 4 & -4\end{pmatrix}$, yang
singular dan semidefinit negatif: jadi ujinya membisu; lalu kajian
berarahnya ($g(x,x) = 2x^4 > 0$ dan $g(x,-x) = 2x^4 - 8x^2 < 0$ untuk yang kecil)
menunjukkan titik bergaya pelana --- jadi tanpa ekstremum.
\end{solution}

\begin{solution}{exo:b2:diffcalc:4}
Turunkan $t \mapsto f(tx)$ di $t = 1$: menurut aturan rantai ia bernilai
$\langle \nabla f(x), x\rangle$; sedangkan berkat kehomogenannya fungsi yang sama
adalah $t^pf(x)$, yang berturunan $pf(x)$ di $t = 1$: itulah kesamaan Euler.

Konversnya: tetapkan $x \neq 0$ lalu misalkan $\varphi(t) = f(tx) - t^p f(x)$
pada $t > 0$. Maka $\varphi'(t) = \langle\nabla f(tx), x\rangle -
pt^{p-1}f(x) = \frac1t\bigl(\langle \nabla f(tx), tx\rangle -
p\,t^pf(x)\bigr)$. Adapun hipotesisnya --- yakni kesamaan Euler di titik
$tx$ --- menilai kurungnya sebagai $p\,f(tx) - p\,t^pf(x) =
p\,\varphi(t)$. Jadi $\varphi' = \frac{p}{t}\varphi$ dengan $\varphi(1)
= 0$: dan penyelesaian tunggal persamaan \omterm{def:b2:diffcalc:differential}{diferensial} linearnya adalah $\varphi \equiv 0$ (berkat
ketunggalan Tahun ke-1), yakni $f(tx) = t^pf(x)$.
\end{solution}

\begin{solution}{exo:b2:diffcalc:5}
Menguraikan $f(x + h) - f(x) = \langle Ax - b, h\rangle +
\frac12\langle Ah, h\rangle$ (berkat kesimetrikan $A$) memberi $\nabla f(x) = Ax -
b$ dan $H_f = A$ di mana-mana. Lalu $f$ cembung bila dan hanya bila $A$ semidefinit
positif (lewat uji Hessiannya, yang global di sini karena $H$ tetap:
sebab rumus Taylor orde duanya eksak). Bila $A$ definit
positif: maka satu-satunya titik kritisnya adalah $x^* = A^{-1}b$, dan $f(x^*
+ h) - f(x^*) = \frac12\langle Ah, h\rangle \geq
\frac{\lambda_{\min}}{2}\norm h^2 > 0$ untuk $h \neq 0$: jadi minimum
global yang tegas.
\end{solution}

\begin{solution}{exo:b2:diffcalc:6}
Karena $\nabla g(a) \neq 0$, satu turunan parsialnya, katakanlah $\frac{\partial
g}{\partial y}(a) \neq 0$: jadi teorema fungsi implisit
(yakni pendamping \cref{thm:b2:diffcalc:inverse}) memparameterkan $\{g =
0\}$ di dekat $a = (a_1, a_2)$ sebagai $\gamma(t) = (t, y(t))$ dengan $y$
bersifat $C^1$ dan $y'(t) = -\frac{\partial_x g}{\partial_y g}(\gamma(t))$
(turunkan $g(t, y(t)) = 0$). Lalu fungsi satu peubah $t
\mapsto f(\gamma(t))$ punya ekstremum lokal di $t = a_1$:
\[
0 = \frac{\dd}{\dd t}f(\gamma(t))\Big|_{a_1}
= \partial_x f(a) + \partial_y f(a)\,y'(a_1)
= \partial_x f(a) - \partial_yf(a)\frac{\partial_x
g(a)}{\partial_y g(a)} :
\]
jadi vektor $\nabla f(a)$ dan $\nabla g(a)$ berkoordinat
sebanding: sehingga $\nabla f(a) = \lambda \nabla g(a)$ dengan $\lambda =
\frac{\partial_y f(a)}{\partial_y g(a)}$.

Penerapannya: pada lingkarannya, $\nabla(xy) = (y, x)$ yang sejajar dengan
$(2x, 2y)$ memaksa $y^2 = x^2$; lalu dengan kendalanya, calonnya
adalah $\pm\bigl(\tfrac{1}{\sqrt2}, \tfrac{1}{\sqrt2}\bigr)$ (bernilai
$\frac12$) dan $\pm\bigl(\tfrac{1}{\sqrt2},
-\tfrac{1}{\sqrt2}\bigr)$ (bernilai $-\frac12$): jadi maksimumnya $\frac12$ dan minimumnya
$-\frac12$ (yang tercapai, sebab lingkarannya \omterm{def:b2:metric:compact}{kompak}).
\end{solution}

\begin{solution}{exo:b2:diffcalc:7}
$\nabla f = (2x - y + 1,\; 2y - x - 1) = 0$: setelah diselesaikan, $x =
-\frac13$ dan $y = \frac13$. Hessiannya $\begin{pmatrix} 2 & -1\\ -1 &
2\end{pmatrix}$, yang definit positif: jadi minimum global bagi $f$ yang
kuadratik itu, bernilai $f\bigl(-\frac13, \frac13\bigr) = -\frac13$.

Pada segitiga $T$: titik kritis dalamnya $(-\frac13,
\frac13) \notin T$ (sebab $x$-nya negatif). Rusuknya: pada $y = 0$ dengan $x \in
\intcc{0}{1}$: $f = x^2 + x$, yang naik, jadi ekstremnya $0$ dan $2$. Pada
$x = 0$: $f = y^2 - y$, dengan minimum $-\frac14$ di $y = \frac12$, dan bernilai
$0$ serta $0$ di ujungnya. Pada $x + y = 1$: sulihkan $y = 1 - x$, jadi
$f = x^2 + (1-x)^2 - x(1-x) + x - (1-x) = 3x^2 - x$; lalu pada
$\intcc{0}{1}$: minimumnya $-\frac{1}{12}$ di $x = \frac16$, dengan nilai
$0$ (di $x=0$) dan $2$ (di $x=1$). Titik sudutnya: $f(0,0) = 0$, $f(1,0)
= 2$, $f(0,1) = 0$. Jadi secara global pada $T$: minimumnya $-\frac14$ di $(0,
\frac12)$, dan maksimumnya $2$ di $(1, 0)$.
\end{solution}

\begin{solution}{exo:b2:diffcalc:8}
$J_f = \begin{pmatrix} 1 & 2y\\ 2x & 1\end{pmatrix}$ dengan $\det J_f =
1 - 4xy$. Di $0$: $\det = 1 \neq 0$, jadi difeomorfisma lokal
(\cref{thm:b2:diffcalc:inverse}), dengan
\[
\dd(f^{-1})_{(0,0)} = (J_f(0))^{-1} = I_2 .
\]
Adapun keterbalikannya pada sebuah bola: kita menuntut $4\abs{xy} < 1$ di seluruhnya; dan pada
$\norm{(x,y)}_2 < r$ berlaku $\abs{xy} \leq \frac{x^2 + y^2}{2} <
\frac{r^2}{2}$, jadi $r = \frac{1}{\sqrt2}$ berhasil; dan itulah yang terbesar:
sebab di $(x, y) = \bigl(\tfrac12, \tfrac12\bigr)$, yang bernorma
$\frac{1}{\sqrt2}$, berlaku $\det J_f = 0$.
\end{solution}

\begin{solution}{exo:b2:diffcalc:9}
$f(t) = (\cos t, \sin t)$: $f(0) = f(2\pi) = (1, 0)$, namun $f'(t) =
(-\sin t, \cos t)$ bernorma $1$ dan tak pernah nol: jadi tak ada titik tempat
turunannya lenyap --- sehingga Rolle tak punya analogi vektor. Adapun ketaksamaan
nilai rerata berlaku dengan nyaman: $\norm{f(2\pi) - f(0)} = 0 \leq
2\pi \cdot \sup\norm{f'} = 2\pi$.
\end{solution}

\begin{solution}{exo:b2:diffcalc:10}
$f(x+h) - f(x) = 2\langle x, h\rangle + \norm h^2$: jadi $\dd f_x =
2\langle x, \cdot\rangle$, $\nabla f(x) = 2x$, dan Hessiannya $2I$
(yang tetap). Untuk $g$:
\[
g(x + h) - g(x) = \langle Ax, h\rangle + \langle Ah, x\rangle
+ \langle Ah, h\rangle
= \bigl\langle (A + A^{\mathsf T})x,\ h\bigr\rangle +
O(\norm h^2),
\]
jadi $\nabla g(x) = (A + A^{\mathsf T})x$ dan $H_g = A +
A^{\mathsf T}$, yang tetap. Menurut \cref{exo:b2:diffcalc:11}, $g$ bersifat
cembung bila dan hanya bila $A + A^{\mathsf T}$ semidefinit positif ---
jadi hanya bagian \omterm{def:b2:quadratic:adjoint}{simetrik} $A$ yang berperan, sebab memang $g(x) =
\langle \frac{A + A^{\mathsf T}}2 x, x\rangle$.
\end{solution}

\begin{solution}{exo:b2:diffcalc:11}
Fungsi $f$ cembung bila dan hanya bila pembatasannya ke setiap ruas cembung,
yakni bila setiap $\varphi(t) = f(a + tv)$ cembung. Menurut aturan
rantai, $\varphi''(t) = \langle H_{a+tv}\,v,\ v\rangle$.

Bila semua Hessiannya semidefinit positif: maka $\varphi'' \geq 0$,
jadi tiap $\varphi$ cembung (lewat jilid Tahun ke-1) sehingga $f$ cembung.
Sebaliknya bila $f$ cembung, tiap $\varphi$ cembung, jadi
$\varphi''(0) \geq 0$: sehingga $\langle H_a v, v\rangle \geq 0$ untuk
setiap $a$ dan setiap arah $v$, yakni semua Hessiannya semidefinit
positif.
\end{solution}

\begin{solution}{exo:b2:diffcalc:12}
\begin{enumerate}
  \item Menurut ketaksamaan nilai rerata
        (\cref{thm:b2:diffcalc:mvi}) yang diterapkan pada $g$:
        $\norm{g(x) - g(y)} \leq k\norm{x-y}$, jadi
        \[
        \norm{f(x) - f(y)} \geq \norm{x - y} - \norm{g(x) -
        g(y)} \geq (1 - k)\norm{x - y} :
        \]
        jadi $f$ bersifat injektif dan $f^{-1}$ (yang terdefinisi pada petanya) bersifat
        \omterm{def:b2:metric:continuity}{Lipschitz} berkonstanta $\frac{1}{1-k}$.
  \item Tetapkan $y$; maka $T(x) = y - g(x)$ memenuhi $\norm{T(x) -
        T(x')} = \norm{g(x') - g(x)} \leq k\norm{x - x'}$: yakni sebuah
        kontraksi pada ruang \omterm{def:b2:metric:complete}{lengkap} $\R^n$. Lalu Banach
        (\cref{thm:b2:metric:banach}) memberi titik tetap $x^*
        = y - g(x^*)$, yakni $f(x^*) = y$: jadi $f$ surjektif.
  \item Jadi $f$ merupakan bijeksi $\R^n$ dengan balikan yang
        \omterm{def:b2:metric:continuity}{Lipschitz} berkonstanta $\frac{1}{1-k}$: yakni sebuah teorema balikan
        global, tempat kekecilan $\dd g$ di mana-mana
        menggantikan hipotesis keterbalikan lokal pada
        \cref{thm:b2:diffcalc:inverse}.
\end{enumerate}
\end{solution}

\begin{solution}{pb:b2:diffcalc:1}
\textbf{1.} Berlaku $\norm{ABx} \leq \vertiii A\norm{Bx} \leq
\vertiii A\vertiii B\norm x$: jadi ambillah sup atas $\norm x = 1$.
Adapun hasil kali matriks, $\det$ dan $\operatorname{tr}$ merupakan fungsi
polinomial atas entrinya, sehingga \omterm{def:b2:metric:continuity}{kontinu} (sebab $\mathcal M_n(\R)
\simeq \R^{n^2}$ dan semua \omterm{def:b2:nvs:norm}{normanya} setara:
\cref{thm:b2:nvs:finitedim}).

\textbf{2.} Berlaku $\sum\vertiii{X^k} \leq \sum\vertiii X^k <
\infty$: jadi deretnya konvergen \omterm{def:b2:series:def}{mutlak}, sehingga konvergen
(\cref{thm:b2:nvs:absoluteconvergence}). Lalu dari $(I -
X)\sum_{k\leq N}X^k = I - X^{N+1} \to I$: jumlahnya adalah $(I -
X)^{-1}$. \omterm{def:b2:nvs:norm}{Normanya}: $\leq \sum\vertiii X^k = \frac{1}{1 -
\vertiii X}$; dan
\[
(I - X)^{-1} - I - X = \sum_{k\geq2}X^k = X^2(I - X)^{-1},
\qquad
\vertiii{X^2(I-X)^{-1}} \leq
\frac{\vertiii X^2}{1 - \vertiii X} = O(\vertiii X^2).
\]

\textbf{3.} Berlaku $A + H = A(I + A^{-1}H)$ dengan $\vertiii{A^{-1}H}
\leq \vertiii{A^{-1}}\vertiii H < 1$: jadi terbalikkan menurut pertanyaan 2,
sehingga bola \omterm{def:b2:metric:topology}{terbuka} berjari-jari $\vertiii{A^{-1}}^{-1}$ di sekitar $A$ terletak
di $GL_n(\R)$. Kepadatannya: $\det(A + \varepsilon I)$ adalah polinomial
berderajat $n$ dalam $\varepsilon$ dengan koefisien utama $1$: jadi ia
punya berhingga banyak akar, sehingga ada $\varepsilon_k \to 0$
dengan $A + \varepsilon_kI$ terbalikkan dan konvergen ke $A$.

\textbf{4.} Untuk $\vertiii H < \frac{1}{2\vertiii{A^{-1}}}$:
\[
(A + H)^{-1} - A^{-1}
= \bigl[(I + A^{-1}H)^{-1} - I\bigr]A^{-1},
\]
yang bernorma paling banyak $\frac{\vertiii{A^{-1}H}}{1 -
\vertiii{A^{-1}H}}\,\vertiii{A^{-1}} \leq
2\vertiii{A^{-1}}^2\vertiii H \to 0$: jadi $\Phi$ \omterm{def:b2:metric:continuity}{kontinu} di
setiap $A \in GL_n(\R)$.

\textbf{5.} Berlaku $(A+H)^2 = A^2 + AH + HA + H^2$: jadi pemetaan $H
\mapsto AH + HA$ bersifat linear dan galatnya $H^2$ bernilai
$O(\vertiii H^2)$. Menguraikan $(A + H)^k$ lalu menyortirnya menurut
banyaknya faktor $H$: bagian linearnya adalah
$\sum_{i=0}^{k-1}A^iHA^{k-1-i}$, sedangkan suku dengan $\geq 2$
faktor $H$ dibatasi oleh sebanyak $\binom k2$ hasil kali bernorma
berskala $\leq \vertiii A^{k-2}\vertiii H^2$: jadi $O(\vertiii H^2)$.
Kita tak dapat meruntuhkan jumlahnya menjadi $kA^{k-1}H$ karena $H$ dan $A$
tak harus komut --- jadi jumlahnya adalah turunan tak komutatif
yang benar.

\textbf{6.} Untuk $H$ yang kecil:
\[
(A+H)^{-1} = (I + A^{-1}H)^{-1}A^{-1}
= \bigl(I - A^{-1}H + O(\vertiii H^2)\bigr)A^{-1}
= A^{-1} - A^{-1}HA^{-1} + O(\vertiii H^2) :
\]
jadi $\dd\Phi_A(H) = -A^{-1}HA^{-1}$, yang linear dalam $H$. Untuk $n = 1$:
$\dd(1/a)(h) = -h/a^2$, yakni turunan yang sudah dikenal.

\textbf{7.} Aturan rantai sepanjang kurvanya memberi
$\bigl(A(t)^{-1}\bigr)' = \dd\Phi_{A(t)}(A'(t)) =
-A(t)^{-1}A'(t)A(t)^{-1}$. Di $A(t) = I + tB$ dengan $t = 0$: $(I +
tB)^{-1} = I - tB + O(t^2)$.

\textbf{8.} Menurut pertanyaan 5 dan keinvarianan \omterm{def:b2:structures:generated}{siklik} tracenya:
\[
\dd f_A(H) = \operatorname{tr}\Bigl(\sum_{i=0}^{k-1}
A^iHA^{k-1-i}\Bigr) = k\operatorname{tr}\bigl(A^{k-1}H\bigr) .
\]
Terhadap hasil kali Frobeniusnya, syarat $\dd f_A(H) =
\operatorname{tr}\bigl((\nabla f)^{\mathsf T}H\bigr)$ menuntut
$(\nabla f)^{\mathsf T} = kA^{k-1}$: jadi $\nabla f(A) =
k\,(A^{k-1})^{\mathsf T}$.

\textbf{9.} Berlaku $f(A + H) - f(A) = 2\operatorname{tr}
(A^{\mathsf T}H) + \operatorname{tr}(H^{\mathsf T}H)$: jadi
\omterm{def:b2:diffcalc:differential}{diferensialnya} adalah $H \mapsto 2\operatorname{tr}(A^{\mathsf T}H) =
2\langle A, H\rangle$, sehingga $\nabla f(A) = 2A$; sedangkan suku
orde duanya persis $\norm H_F^2$: jadi Hessiannya dua kali
\omterm{def:b2:quadratic:def}{bentuk kuadratik} identitasnya, yang definit positif dan tetap, sehingga
$\norm\cdot_F^2$ bersifat cembung tegas (sebab rumus Taylornya
eksak di sini).

\textbf{10.} Berkat kemultilinearan pada kolomnya, $\det(I + H) =
\sum_{S\subseteq\{1,\dots,n\}}\det(M_S)$ dengan $M_S$ berkolom
$h_j$ untuk $j \in S$ dan $e_j$ selebihnya. Adapun $S = \varnothing$
memberi $1$; lalu $S = \{j\}$ memberi \omterm{def:b2:linalg:det}{determinan} $I$ dengan kolom
$j$-nya diganti $h_j$, yakni entri ke-$j$-nya $h_{jj}$,
yang berjumlah $\operatorname{tr}H$; sedangkan tiap suku dengan $\abs S \geq 2$
adalah \omterm{def:b2:linalg:det}{determinan} dengan sekurangnya dua kolom berukuran
$O(\vertiii H)$, sehingga $O(\vertiii H^2)$ (sebab pemetaan multilinear pada ruang
berdimensi hingga bersifat terbatas), dan banyaknya berhingga.
Jadi $\det(I + H) = 1 + \operatorname{tr}H +
O(\vertiii H^2)$: sehingga $\dd(\det)_I = \operatorname{tr}$.

\textbf{11.} Berlaku $\det(A + H) = \det A\,\det(I + A^{-1}H) =
\det A\,\bigl(1 + \operatorname{tr}(A^{-1}H) +
O(\vertiii H^2)\bigr)$: jadi \omterm{def:b2:diffcalc:differential}{diferensialnya} adalah $H \mapsto
\det(A)\operatorname{tr}(A^{-1}H)$.

\textbf{12.} Penguraian Laplace sepanjang baris $i$: $\det A =
\sum_j a_{ij}C_{ij}$, dan kofaktor $C_{ij}$ tak melibatkan
baris $i$: jadi $\frac{\partial\det}{\partial a_{ij}} = C_{ij}$. Karena itu
\[
\dd(\det)_A(H) = \sum_{i,j}C_{ij}h_{ij}
= \operatorname{tr}\bigl(\operatorname{com}(A)^{\mathsf T}
H\bigr)
= \operatorname{tr}\bigl(\operatorname{adj}(A)H\bigr),
\qquad
\nabla(\det)(A) = \operatorname{com}(A) .
\]
Untuk $A$ yang terbalikkan, $\operatorname{adj}A = \det(A)A^{-1}$
memulihkan pertanyaan 11. Sepanjang kurva $C^1$, aturan rantainya berbunyi
$(\det A(t))' = \operatorname{tr}(\operatorname{adj}
(A(t))\,A'(t))$: itulah rumus Jacobi.

\textbf{13.} Dengan $A' = MA$ dan
$\operatorname{adj}(A)A = \det(A)I$:
\[
(\det A)' = \operatorname{tr}\bigl(\operatorname{adj}(A)MA
\bigr)
= \operatorname{tr}\bigl(A\operatorname{adj}(A)M\bigr)
= \det A\;\operatorname{tr}M
\]
(berkat kesiklikannya; sebab $A\operatorname{adj}A = \det(A) I$ juga). Adapun
persamaan \omterm{def:b2:diffcalc:differential}{diferensial} linear skalar $y' = \operatorname{tr}(M(t))\,y$ punya
penyelesaian tunggal $y(t) = y(0)\exp\bigl(\int_0^t
\operatorname{tr}M\bigr)$ (lewat Tahun ke-1): itulah rumus Liouville.

\textbf{14.} Di $A \in SL_n(\R)$, ambillah $H = \frac1nA$: maka
$\dd(\det)_A\bigl(\tfrac1nA\bigr) = \frac1n\det(A)
\operatorname{tr}(A^{-1}A) = \frac1n\cdot1\cdot n = 1 \neq 0$:
jadi \omterm{def:b2:diffcalc:differential}{diferensialnya} merupakan bentuk linear yang tak nol, sehingga surjektif
ke $\R$ di setiap titik himpunan arasnya: jadi $SL_n(\R)$ merupakan
himpunan aras yang mulus (berdimensi $n^2 - 1$).

\textbf{15.} Berlaku $\sum_k\vertiii{A^k/k!} \leq
\sum\vertiii A^k/k! = \eu^{\vertiii A}$: jadi kekonvergenan \omterm{def:b2:series:def}{mutlaknya}
(lewat hujah kelengkapan pertanyaan 2), dengan \omterm{def:b2:funcseq:series}{kekonvergenan normal}
pada setiap bola $\vertiii A \leq R$ (sebab batasnya $R^k/k!$ tak bergantung
pada $A$). Lalu tiap jumlah parsialnya \omterm{def:b2:metric:continuity}{kontinu} (sebab polinomial); dan
limit \omterm{def:b2:funcseq:def}{seragamnya} pada bola bersifat \omterm{def:b2:metric:continuity}{kontinu}: jadi $A \mapsto \eu^A$ bersifat
\omterm{def:b2:metric:continuity}{kontinu}, dengan $\vertiii{\eu^A} \leq \eu^{\vertiii A}$.

\textbf{16.} Kedua deretnya konvergen \omterm{def:b2:series:def}{mutlak}, jadi \omterm{thm:b2:series:fubini}{hasil kali
Cauchynya} sah (\cref{thm:b2:series:fubini}):
\[
\eu^A\eu^B = \sum_{n\geq0}\frac{1}{n!}\sum_{k=0}^n\binom
nkA^kB^{n-k}
= \sum_{n\geq0}\frac{(A+B)^n}{n!} = \eu^{A+B},
\]
dengan kesamaan binomialnya menuntut $AB = BA$. Lalu dengan $B = -A$:
$\eu^A\eu^{-A} = \eu^0 = I$: jadi setiap $\eu^A \in GL_n(\R)$.

\textbf{17.} Deret $\sum t^kA^k/k!$ dan deret turunannya
$\sum t^{k-1}A^k/(k-1)! = A\sum t^{k-1}A^{k-1}/(k-1)!$ konvergen
normal pada setiap ruas $\abs t \leq T$ (dengan batas
$T^k\vertiii A^k/k!$): jadi teorema penurunan bagi deret
(\cref{thm:b2:funcseq:seriestransfer}, yang diterapkan entri demi entri)
memberi $\frac{\dd}{\dd t}\eu^{tA} = A\eu^{tA}$; sedangkan memfaktorkan $A$
di kanannya justru memberi $\eu^{tA}A$. Setelah diiterasikan: jadi $C^\infty$.

\textbf{18.} Kurva $A(t) = \eu^{tA}$ memenuhi $A'(t) = A\,A(t)$:
jadi rumus Liouville (pertanyaan 13) dengan $M = A$ yang tetap memberi
$\det\eu^{tA} = \eu^{t\operatorname{tr}A}$ (yang bernilai $1$ di $t =
0$); lalu di $t = 1$: $\det\eu^A = \eu^{\operatorname{tr}A}$.
Pemeriksaannya: kasus $n = 1$ adalah eksponensialnya sendiri; lalu $A$ yang nilpoten punya
$\operatorname{tr}A = 0$ dengan $\eu^A$ unipoten berdeterminan
$1$; dan $\operatorname{tr}A = 0$ memberi $\det\eu^A = 1$: jadi
matriks bertrace nol dikirim ke dalam $SL_n(\R)$.

\textbf{19.} Berlaku $\eu^H - I - H = \sum_{k\geq2}H^k/k!$, yang bernorma
$\leq \vertiii H^2\eu^{\vertiii H} = O(\vertiii H^2)$: jadi
$\dd(\exp)_0 = \mathrm{id}$, yang terbalikkan. Lebih jauh $\exp$ bersifat
$C^1$: sebab menurut pertanyaan 5, calon \omterm{def:b2:diffcalc:differential}{diferensialnya} $H \mapsto
\sum_k\frac1{k!}\sum_iA^iHA^{k-1-i}$ merupakan deret pemetaan linear yang
konvergen normal dan bergantung secara \omterm{def:b2:metric:continuity}{kontinu} pada $A$ (dengan batas
$\vertiii A^{k-1}/(k-1)!$ pada bola), jadi turunan parsialnya ada dan
\omterm{def:b2:metric:continuity}{kontinu} (lewat \cref{thm:b2:diffcalc:c1} dan teorema pemindahan
deretnya). Lalu teorema fungsi balikan
(\cref{thm:b2:diffcalc:inverse}) berlaku di $0$: jadi $\exp$ merupakan
difeomorfisma $C^1$ dari sebuah persekitaran $0$ ke sebuah
persekitaran $I$ --- sehingga matriks di dekat $I$ punya logaritma.

\textbf{20.} Untuk $S = PDP^{\mathsf T}$ yang \omterm{def:b2:quadratic:adjoint}{simetrik} (lewat teorema
spektral): $\eu^S = P\eu^DP^{\mathsf T}$ bersifat \omterm{def:b2:quadratic:adjoint}{simetrik} dengan
\omterm{def:b2:reduction:eigen}{nilai eigen} $\eu^{\lambda_i} > 0$: jadi definit positif.
Kesurjektifannya: sebab $Q =
P\operatorname{diag}(\mu_i)P^{\mathsf T}$ yang definit positif (dengan $\mu_i > 0$) sama dengan
$\eu^S$ untuk $S = P\operatorname{diag}(\ln\mu_i)P^{\mathsf T}$.
Keinjektifannya: sebab $\eu^S$ menentukan \omterm{def:b2:reduction:eigen}{ruang eigennya}, yang
persis merupakan ruang eigen $S$ (sebab pada tiap ruang eigen $S$ bagi $\lambda$,
$\eu^S$ bekerja sebagai $\eu^\lambda$; dan $\lambda$ yang berbeda memberi
$\eu^\lambda$ yang berbeda), lalu mengambil $\ln$ \omterm{def:b2:reduction:eigen}{nilai eigennya}
memulihkan $S$. Jadi $\exp$ merupakan bijeksi dari matriks \omterm{def:b2:quadratic:adjoint}{simetrik}
ke matriks yang definit positif.

\textbf{21.} Berlaku $F(A + H) = A^{\mathsf T}A + A^{\mathsf T}H +
H^{\mathsf T}A + H^{\mathsf T}H$: jadi $\dd F_A(H) =
A^{\mathsf T}H + H^{\mathsf T}A$ (yang bernilai di $S_n$, dengan galat
$O(\vertiii H^2)$). Lalu di $A \in O_n$ dan untuk $S \in S_n$,
pilihan $H = \frac12AS$ memberi
\[
A^{\mathsf T}\cdot\tfrac12AS + \tfrac12(AS)^{\mathsf T}A
= \tfrac12 S + \tfrac12 S^{\mathsf T} = S :
\]
jadi surjektif. Sehingga $O_n = F^{-1}(I)$ merupakan himpunan aras yang mulus,
berdimensi $n^2 - \dim S_n = \frac{n(n-1)}2$.

\textbf{22.} Menurunkan $A(t)^{\mathsf T}A(t) = I$ di $t
= 0$ (dengan $A(0) = I$) memberi $A'(0)^{\mathsf T} + A'(0) = 0$: jadi
antisimetrik. Sebaliknya, untuk $K^{\mathsf T} = -K$:
$(\eu^{tK})^{\mathsf T}\eu^{tK} = \eu^{tK^{\mathsf T}}
\eu^{tK} = \eu^{-tK}\eu^{tK} = I$ (transposkan deretnya
suku demi suku; sebab eksponennya komut): jadi kurvanya tinggal di $O_n$,
dengan kecepatan $K$ di $t = 0$. Jadi ruang singgungnya di $I$ $=$
matriks antisimetriknya, yang berdimensi seperti diharapkan
$\frac{n(n-1)}2$.

\textbf{23.} Berlaku $\operatorname{tr}K = 0$ untuk $K$ yang antisimetrik,
jadi $\det\eu^K = \eu^0 = 1$ (pertanyaan 18): sehingga eksponensialnya
mendarat di $SO_n$. Untuk $n = 2$: $J^2 = -I$, jadi $J^{2m} =
(-1)^mI$ dan $J^{2m+1} = (-1)^mJ$, sehingga
\[
\eu^{\theta J}
= \Bigl(\sum_m\frac{(-1)^m\theta^{2m}}{(2m)!}\Bigr)I
+ \Bigl(\sum_m\frac{(-1)^m\theta^{2m+1}}{(2m+1)!}\Bigr)J
= \cos\theta\,I + \sin\theta\,J ,
\]
yakni perputaran sebesar $\theta$: jadi deret sinus dan kosinusnya hidup
di dalam eksponensial matriksnya.

\textbf{24.} Untuk $A, B$ yang terbalikkan: $\operatorname{adj}(AB) =
\det(AB)(AB)^{-1} = \det(B)\det(A)B^{-1}A^{-1} =
\operatorname{adj}(B)\operatorname{adj}(A)$. Adapun kedua ruas
kesamaannya merupakan pemetaan polinomial (sehingga \omterm{def:b2:metric:continuity}{kontinu}) atas entri
$(A, B)$; dan keduanya sepakat pada himpunan bagian padat $GL_n\times GL_n$
milik $\mathcal M_n\times\mathcal M_n$ (lewat pertanyaan 3: hampirilah
tiap faktornya), jadi keduanya sepakat di mana-mana.

\textbf{25.} (i) Bagian I berjalan atas kelengkapan ruang bernorma
berdimensi hingga (sebab deret yang konvergen \omterm{def:b2:series:def}{mutlak} akan konvergen), Bagian IV
atas hal yang sama ditambah teorema spektral untuk pertanyaan 20, dan logaritma
lokal Bagian V atas teorema fungsi balikan. (ii) Bab
persamaan \omterm{def:b2:diffcalc:differential}{diferensial} menyandari rumus Liouville
(pertanyaan 13) bagi Wronskian sistem linearnya, dan
$\frac{\dd}{\dd t}\eu^{tA} = A\eu^{tA}$ (pertanyaan 17), yang tak lain
pernyataan bahwa $\eu^{tA}$ menyelesaikan $X' = AX$. (iii)
Kesamaan $\dd(\det)_I = \operatorname{tr}$ mengatakan bahwa trace adalah
laju infinitesimal perubahan volumenya, sedangkan $\det\eu^A =
\eu^{\operatorname{tr}A}$ mengintegralkan pernyataan itu secara global.
(iv) Jilid Tahun ke-3 menamai strukturnya: $O_n$ dan
$SL_n(\R)$ adalah grup Lie, ruang singgungnya di $I$
(yakni matriks antisimetrik dan bertrace nol) adalah aljabar Lie, dan
$\exp$ adalah jembatan di antara keduanya.

\end{solution}
